\begin{lemma}
\label{lemma-lift-complex}
Let $\mathcal{C}$ be a site. Let $\mathcal{O} \to \mathcal{O}_0$
be a surjection of sheaves of rings. Assume given the following data
\begin{enumerate}
\item flat $\mathcal{O}$-modules $\mathcal{G}^n$,
\item maps of $\mathcal{O}$-modules $\mathcal{G}^n \to \mathcal{G}^{n + 1}$,
\item a complex $\mathcal{K}_0^\bullet$ of $\mathcal{O}_0$-modules,
\item maps of $\mathcal{O}$-modules $\mathcal{G}^n \to \mathcal{K}_0^n$
\end{enumerate}
such that
\begin{enumerate}
\item[(a)] $H^n(\mathcal{K}_0^\bullet) = 0$ for $n \gg 0$,
\item[(b)] $\mathcal{G}^n = 0$ for $n \gg 0$,
\item[(c)] with
$\mathcal{G}^n_0 = \mathcal{G}^n \otimes_\mathcal{O} \mathcal{O}_0$
the induced maps determine a complex $\mathcal{G}_0^\bullet$ and a map
of complexes $\mathcal{G}_0^\bullet \to \mathcal{K}_0^\bullet$.
\end{enumerate}
Then there exist
\begin{enumerate}
\item[(\romannumeral1)]
flat $\mathcal{O}$-modules $\mathcal{F}^n$,
\item[(\romannumeral2)]
maps of $\mathcal{O}$-modules $\mathcal{F}^n \to \mathcal{F}^{n + 1}$,
\item[(\romannumeral3)]
maps of $\mathcal{O}$-modules $\mathcal{F}^n \to \mathcal{K}_0^n$,
\item[(\romannumeral4)]
maps of $\mathcal{O}$-modules $\mathcal{G}^n \to \mathcal{F}^n$,
\end{enumerate}
such that $\mathcal{F}^n = 0$ for $n \gg 0$, such that the diagrams
$$
\xymatrix{
\mathcal{G}^n \ar[r] \ar[d] & \mathcal{G}^{n + 1} \ar[d] \\
\mathcal{F}^n \ar[r] & \mathcal{F}^{n + 1}
}
$$
commute for all $n$, such that the composition
$\mathcal{G}^n \to \mathcal{F}^n \to \mathcal{K}_0^n$
is the given map $\mathcal{G}^n \to \mathcal{K}_0^n$, and such that with
$\mathcal{F}^n_0 = \mathcal{F}^n \otimes_\mathcal{O} \mathcal{O}_0$
we obtain a complex $\mathcal{F}_0^\bullet$ and map of complexes
$\mathcal{F}_0^\bullet \to \mathcal{K}_0^\bullet$ which is a
quasi-isomorphism.
\end{lemma}

\begin{proof}
We will prove by descending induction on $e$ that we can find $\mathcal{F}^n$,
$\mathcal{G}^n \to \mathcal{F}^n$, and
$\mathcal{F}^n \to \mathcal{F}^{n + 1}$ for $n \geq e$
fitting into a commutative diagram
$$
\xymatrix{
\ldots \ar[r] &
\mathcal{G}^{e - 1} \ar[r] \ar@/_2pc/[dd] &
\mathcal{G}^e \ar[d] \ar[r] \ar@/_2pc/[dd] &
\mathcal{G}^{e + 1} \ar[d] \ar[r] \ar@/_2pc/[dd]|\hole &
\ldots \\
& &
\mathcal{F}^e \ar[d] \ar[r] &
\mathcal{F}^{e + 1} \ar[d] \ar[r] & \ldots \\
\ldots  \ar[r] &
\mathcal{K}_0^{e - 1} \ar[r] &
\mathcal{K}_0^e \ar[r] &
\mathcal{K}_0^{e + 1} \ar[r] & \ldots
}
$$
such that $\mathcal{F}_0^\bullet$ is a complex,
the induced map $\mathcal{F}_0^\bullet \to \mathcal{K}_0^\bullet$
induces an isomorphism on $H^n$ for $n > e$ and a surjection
for $n = e$. For $e \gg 0$ this is true because we can take
$\mathcal{F}^n = 0$ for $n \geq e$ in that case by assumptions
(a) and (b).

\medskip\noindent
Induction step. We have to construct $\mathcal{F}^{e - 1}$
and the maps $\mathcal{G}^{e - 1} \to \mathcal{F}^{e - 1}$,
$\mathcal{F}^{e - 1} \to \mathcal{F}^e$, and
$\mathcal{F}^{e - 1} \to \mathcal{K}_0^{e - 1}$.
We will choose $\mathcal{F}^{e - 1} = A \oplus B \oplus C$
as a direct sum of three pieces.

\medskip\noindent
For the first we take $A = \mathcal{G}^{e - 1}$ and we choose our map
$\mathcal{G}^{e - 1} \to \mathcal{F}^{e - 1}$ to be the inclusion of
the first summand. The maps $A \to \mathcal{K}^{e - 1}_0$
and $A \to \mathcal{F}^e$ will be the obvious ones.

\medskip\noindent
To choose $B$ we consider the surjection (by induction hypothesis)
$$
\gamma :
\Ker(\mathcal{F}^e_0 \to \mathcal{F}^{e + 1}_0)
\longrightarrow
\Ker(\mathcal{K}^e_0 \to \mathcal{K}^{e + 1}_0)/
\Im(\mathcal{K}^{e - 1}_0 \to \mathcal{K}^e_0)
$$
We can choose a set $I$, for each $i \in I$
an object $U_i$ of $\mathcal{C}$, and sections
$s_i \in \mathcal{F}^e(U_i)$, $t_i \in \mathcal{K}^{e - 1}_0(U_i)$
such that
\begin{enumerate}
\item $s_i$ maps to a section of $\Ker(\gamma) \subset
\Ker(\mathcal{F}^e_0 \to \mathcal{F}^{e + 1}_0)$,
\item $s_i$ and $t_i$ map to the same section of
$\mathcal{K}^e_0$,
\item the sections $s_i$ generate $\Ker(\gamma)$ as an $\mathcal{O}_0$-module.
\end{enumerate}
We omit giving the full justification for this;
one uses that $\mathcal{F}^e \to \mathcal{F}^e_0$
is a surjective maps of sheaves of sets. Then we set
to put
$$
B = \bigoplus\nolimits_{i \in I} j_{U_i!}\mathcal{O}_{U_i}
$$
and define the maps $B \to \mathcal{F}^e$ and $B \to \mathcal{K}_0^{e - 1}$
by using $s_i$ and $t_i$ to determine where to send the summand
$j_{U_i!}\mathcal{O}_{U_i}$.

\medskip\noindent
With $\mathcal{F}^{e - 1} = A \oplus B$ and maps as above,
this produces a diagram as above for $e - 1$ such that
$\mathcal{F}_0^\bullet \to \mathcal{K}_0^\bullet$
induces an isomorphism on $H^n$ for $n \geq e$.
To get the map to be surjective on $H^{e - 1}$ we choose
the summand $C$ as follows.
Choose a set $J$, for each $j \in J$ an object $U_j$ of $\mathcal{C}$
and a section $t_j$ of $\Ker(\mathcal{K}^{e - 1}_0 \to \mathcal{K}^e_0)$
over $U_j$ such that these sections generate this kernel over
$\mathcal{O}_0$. Then we put
$$
C = \bigoplus\nolimits_{j \in J} j_{U_j!}\mathcal{O}_{U_j}
$$
and the zero map $C \to \mathcal{F}^e$ and the map
$C \to \mathcal{K}_0^{e - 1}$ by using $s_j$ to determine where to the summand
$j_{U_j!}\mathcal{O}_{U_j}$. This finishes the induction step
by taking $\mathcal{F}^{e - 1} = A \oplus B \oplus C$ and
maps as indicated.
\end{proof}
